\documentclass[12pt,letterpaper]{article}

\usepackage[utf8]{inputenc}
\usepackage[T1]{fontenc}
\usepackage{geometry}
\usepackage{listings}
\usepackage{xcolor}
\usepackage{fancyhdr}
\usepackage{titlesec}
\usepackage{parskip}
\usepackage{booktabs}
\usepackage{array}
\usepackage{mdframed}
\usepackage{hyperref}
\usepackage{enumitem}
\usepackage{tocloft}
\usepackage{tabularx}
\usepackage{longtable}

\geometry{top=1.25in, bottom=1.25in, left=1.25in, right=1.25in}

\definecolor{codebg}{RGB}{245,245,243}
\definecolor{codeborder}{RGB}{200,198,190}
\definecolor{notebg}{RGB}{255,252,235}
\definecolor{noteborder}{RGB}{180,160,80}
\definecolor{warnbg}{RGB}{255,245,245}
\definecolor{warnborder}{RGB}{200,100,100}
\definecolor{stepbg}{RGB}{248,248,248}
\definecolor{stepborder}{RGB}{90,90,90}

\lstset{
  basicstyle=\ttfamily\small,
  backgroundcolor=\color{codebg},
  frame=single,
  framesep=6pt,
  rulecolor=\color{codeborder},
  breaklines=true,
  keepspaces=true,
  columns=flexible,
  xleftmargin=6pt,
  xrightmargin=6pt,
  aboveskip=8pt,
  belowskip=8pt
}

\mdfdefinestyle{notebox}{
  backgroundcolor=notebg, linecolor=noteborder, linewidth=1pt,
  innerleftmargin=12pt, innerrightmargin=12pt,
  innertopmargin=10pt, innerbottommargin=10pt,
  skipabove=10pt, skipbelow=10pt
}
\mdfdefinestyle{warnbox}{
  backgroundcolor=warnbg, linecolor=warnborder, linewidth=1pt,
  innerleftmargin=12pt, innerrightmargin=12pt,
  innertopmargin=10pt, innerbottommargin=10pt,
  skipabove=10pt, skipbelow=10pt
}
\mdfdefinestyle{stepbox}{
  backgroundcolor=stepbg, linecolor=stepborder, linewidth=1pt,
  innerleftmargin=12pt, innerrightmargin=12pt,
  innertopmargin=10pt, innerbottommargin=10pt,
  skipabove=10pt, skipbelow=10pt
}

\pagestyle{fancy}
\fancyhf{}
\fancyhead[L]{\small\textit{AI Document Library --- User Guide}}
\fancyhead[R]{\small\textit{Practical Operation}}
\fancyfoot[C]{\small\thepage}
\renewcommand{\headrulewidth}{0.4pt}

\titleformat{\section}{\large\bfseries}{}{0em}{\thesection\quad}[\vspace{2pt}\hrule\vspace{4pt}]
\titleformat{\subsection}{\normalsize\bfseries}{\thesubsection\quad}{0em}{}
\titleformat{\subsubsection}{\normalsize\itshape}{\thesubsubsection\quad}{0em}{}

\setcounter{secnumdepth}{3}
\setcounter{tocdepth}{3}
\setlist[itemize]{topsep=4pt, itemsep=3pt, parsep=0pt, leftmargin=1.4em}
\setlist[enumerate]{topsep=4pt, itemsep=3pt, parsep=0pt, leftmargin=1.6em}

\hypersetup{colorlinks=true, linkcolor=black, urlcolor=black, bookmarksnumbered=true}
\renewcommand{\cftsecleader}{\cftdotfill{\cftdotsep}}

% ---------------------------------------------------------------
\begin{document}

\begin{titlepage}
\centering
\vspace*{1in}
{\Huge\bfseries AI Document Library}\\[10pt]
{\Large User Guide}\\[8pt]
{\large How to Operate the Library, End to End}\\[20pt]
\hrule height 0.8pt
\vspace{14pt}
{\normalsize
A task-oriented companion to ARCHITECTURE.tex.\\[8pt]
This document explains what to do, in what order, to RESUME a project,\\
run a CHECKPOINT, leave a NOTE, COMMIT routine work, or start a NEW PROJECT.\\
It assumes nothing about which AI vendor or platform you are using.
}
\vspace{14pt}
\hrule height 0.4pt
\vspace{28pt}
{\small Version 1.0 \quad --- \quad \today}
\vfill
{\footnotesize Companion to ARCHITECTURE.tex. See MASTER-PROMPT.md for the\\
authoritative ritual definitions; this guide explains how to act on them.}
\end{titlepage}

\tableofcontents
\newpage

% ---------------------------------------------------------------
\section{Before You Start}

This guide is practical, not philosophical. For the reasoning behind the
five-layer model, see ARCHITECTURE.tex. For the authoritative, word-for-word
ritual definitions, see MASTER-PROMPT.md --- this guide explains how to act
on those rituals, but MASTER-PROMPT.md is always the final word if the two
ever disagree.

\subsection{Document History}

This is version 1.0, written 2026-06-28. There was no prior USER-GUIDE.tex
at this scope: the original USER-GUIDE.tex (2026-04-11) documented a
clipboard copy-paste workflow --- pasting MASTER-PROMPT.md and four RESUME
files into a chat box by hand, once per session. That workflow no longer
exists. RESUME, CHECKPOINT, NOTE, COMMIT, and NEW PROJECT are now file-based
operations carried out by the \texttt{ai-library-ops} skill, and
\texttt{CLAUDE.md} loads MASTER-PROMPT.md automatically in Claude Code,
cloud, and iOS sessions. The original document is archived at
\texttt{projects/ai-library-system/docs/2026-04-11--user-guide-v1-original--claude.tex}
as the historical record of that earlier design. This guide replaces it
entirely rather than amending it, because the underlying workflow changed,
not just the details.

\subsection{What You Need}

\begin{itemize}
\item A clone of the AI-Library git repository.
\item An AI session with file read/write access to that clone --- Claude
      Code (CLI, cloud, or iOS), Cowork, or any platform that can read and
      write files in the repository.
\item Git, configured with \texttt{core.hooksPath} pointed at
      \texttt{code/githooks} (see \nameref{sec:hooks}).
\end{itemize}

\begin{mdframed}[style=notebox]
\textbf{If your platform cannot access files} (a plain chat box with no
file access), Layer 1 still works by hand: paste MASTER-PROMPT.md once,
then paste the four RESUME files when resuming, and paste the checkpoint
blocks back to yourself to save. Every instruction in this guide that
refers to a script or skill running automatically has a manual equivalent
in MASTER-PROMPT.md. Layer 1 never depends on Layer 3 being available.
\end{mdframed}

% ---------------------------------------------------------------
\section{One-Time Setup}
\label{sec:hooks}

Run once per fresh clone. A fresh clone never receives \texttt{.git/hooks/}
automatically, so the portable hooks in \texttt{code/githooks/} have to be
pointed to explicitly:

\begin{lstlisting}
cd AI-Library
git config core.hooksPath code/githooks
chmod +x code/githooks/pre-commit
\end{lstlisting}

If git identity is not already configured in this environment (common on
a fresh cloud or iOS clone):

\begin{lstlisting}
git config user.email "you@example.com"
git config user.name  "Your Name"
\end{lstlisting}

If you are using the \texttt{ai-library-ops} skill (the normal case in
Claude Code), it bootstraps both of the above automatically on its first
operation in a fresh clone. You do not need to do this by hand on every
machine --- only verify it once if something seems to be silently failing.

% ---------------------------------------------------------------
\section{Folder and File Conventions}

\subsection{Where Things Live}

\begin{lstlisting}
AI-Library/
|-- MAP.md                  Index of everything in the library
|-- MASTER-PROMPT.md        The operative system prompt (Layer 2 control)
|-- CLAUDE.md                Imports MASTER-PROMPT.md for Claude Code
|-- ARCHITECTURE.tex        Why the system is built this way
|-- USER-GUIDE.tex          This document
|-- temp/                    Checkpoint staging only
|-- logs/                    Automation run logs
|-- code/                    Deployed scripts and skills
|   `-- githooks/            Portable git hooks
|-- .claude/skills/          Repo-local skills (load in cloud/iOS too)
|-- projects/
|   `-- [project-slug]/
|       |-- THREAD.md             Checkpoint log + inter-checkpoint notes
|       |-- persona.md            Read during RESUME, never preloaded
|       |-- docs/                 Standalone dated documents
|       |-- code/                 Project-local scripts/skills (source)
|       |-- v[NN]--artifact.[ext]
|       |-- v[NN]--context.md
|       `-- v[NN]--instructions.md
`-- inbox/                   Unsorted files awaiting a project
\end{lstlisting}

\subsection{Naming Rules}

\begin{itemize}
\item \textbf{Project triplet files:} \texttt{v[NN]--artifact.[ext]},
      \texttt{v[NN]--context.md}, \texttt{v[NN]--instructions.md}.
      Version numbers are always two digits: \texttt{v01}, \texttt{v02},
      \ldots \texttt{v09}, \texttt{v10}, \texttt{v11}.
\item \textbf{Standalone files} (anything in a \texttt{docs/} folder or
      \texttt{inbox/}): \texttt{YYYY-MM-DD--short-slug--vendor.[ext]}.
\end{itemize}

\subsection{Frontmatter}

Every standalone file (anything named with the dated convention above)
begins with:

\begin{lstlisting}
---
title:
date:
updated:
type: document | code | context
vendor: claude | gpt | gemini | other
model:
tags: []
related: []
---
\end{lstlisting}

The version-triplet files (\texttt{v[NN]--*}) do not use this frontmatter
block --- their structure is already fixed by the checkpoint schema.
Core library files (\texttt{MAP.md}, \texttt{MASTER-PROMPT.md},
\texttt{CLAUDE.md}, \texttt{THREAD.md}, \texttt{persona.md}) are exempt
as well. \LaTeX{} files cannot use a YAML \texttt{---} block without
breaking compilation, so archived or reference \texttt{.tex} files instead
carry the same fields as a \texttt{\%}-commented block before
\texttt{\textbackslash documentclass}.

% ---------------------------------------------------------------
\section{RESUME --- Loading a Project}

Use RESUME to pick up a project you have worked on before. It is the most
common operation in a normal session.

\begin{mdframed}[style=stepbox]
\textbf{What you say:} ``Resume'' (or ``resume a project'', ``load
[project]'', ``switch project'').

\textbf{What happens:}
\begin{enumerate}
\item The AI locates the library root and lists every folder under
      \texttt{projects/}.
\item It asks which project to resume (or proceeds directly if you named
      one).
\item It finds the highest-numbered \texttt{v[NN]--context.md} and
      \texttt{v[NN]--instructions.md} for that project.
\item It reads, in this exact order: \texttt{persona.md}, then
      \texttt{THREAD.md}, then the latest \texttt{context.md}, then the
      latest \texttt{instructions.md}.
\item It states the project, the current version, and a one-sentence
      summary of artifact state, then waits for direction.
\end{enumerate}
\end{mdframed}

\begin{mdframed}[style=notebox]
RESUME reads files from disk. You never need to paste anything --- if an
AI asks you to paste persona.md or THREAD.md content, it does not have
file access to this repository, and you are effectively back on the
Layer 1 manual path.
\end{mdframed}

\subsection{What persona.md Is --- and Is Not}

\texttt{persona.md} is read at RESUME time, as one of the four files
above. It is \emph{not} preloaded into a system prompt and does not
persist automatically between sessions on its own. This was tried early
in the project and abandoned: a system prompt is not designed for
frequent updates, and persona content needs to sit in the same
RESUME-time audit trail as everything else, not in a separate channel
that can drift out of sync with it. If a platform offers a persistent
system-prompt or custom-instructions box, you may paste persona.md there
once as a convenience, but RESUME does not depend on it and will re-read
the file from disk regardless.

% ---------------------------------------------------------------
\section{CHECKPOINT --- Saving Progress}

Use CHECKPOINT to save the current state of a project as a new version.
This is the core ritual of the library.

\subsection{The Four Blocks}

Saying CHECKPOINT produces four blocks, each wrapped in tilde fences with
its label as the first line:

\begin{enumerate}
\item \textbf{ARTIFACT v[NN]} --- the full artifact as raw text. If the
      artifact has not changed since the previous version, the body is
      the single line \texttt{NO CHANGE}; the automation copies and
      relabels the previous artifact file rather than asking you to
      retype it.
\item \textbf{CONTEXT v[NN]} --- PROJECT / DECISIONS / RULED OUT / OPEN /
      STATE. DECISIONS and RULED OUT are append-only: every line from the
      previous context file is copied forward verbatim, then new lines
      are added at the end. Nothing already there is ever removed,
      shortened, or reworded.
\item \textbf{INSTRUCTIONS v[NN]} --- a complete resume packet: GOAL,
      BACKGROUND, ARTIFACT STATE, KEY DECISIONS, OPEN QUESTIONS,
      EXPLICITLY RULED OUT, NEXT TASK, PERSONA, STYLE AND CONSTRAINTS.
\item \textbf{THREAD ENTRY v[NN]} --- the dated log line appended to
      THREAD.md: trigger, artifact state, context state, instructions
      state, key decisions, still-open items.
\end{enumerate}

The version number is always the last THREAD ENTRY number plus one. The
AI determines this by reading THREAD.md, not by asking you.

\subsection{From Blocks to Files}

\begin{mdframed}[style=stepbox]
\begin{enumerate}
\item The four blocks are assembled into one file:
      \texttt{temp/v[NN]-checkpoint.txt}.
\item \textbf{Dry run:} \texttt{python3 code/checkpoint.py [slug] [NN]}.
      This parses the blocks, validates them, and shows exactly what it
      would write --- but writes nothing yet.
\item Read the dry-run output. Confirm the file paths and sizes look
      right.
\item \textbf{Write:} \texttt{python3 code/checkpoint.py [slug] [NN]
      --write}. This writes the versioned triplet, appends the THREAD
      ENTRY to THREAD.md, and updates MAP.md.
\item \textbf{If the artifact is MASTER-PROMPT.md specifically:} copy
      \texttt{v[NN]--artifact.md} to \texttt{MASTER-PROMPT.md} at the
      library root. \texttt{CLAUDE.md} then loads it automatically on
      every future Claude Code, cloud, or iOS session --- no manual
      paste. Only platforms with a separate system-prompt box outside
      this repository (for example, a Cowork project-instructions field)
      need it pasted by hand.
\item \texttt{git add .} and commit. The pre-commit hook validates the
      triplet before allowing the commit through (see
      \nameref{sec:hook-detail}).
\end{enumerate}
\end{mdframed}

If you are using the \texttt{ai-library-ops} skill, saying CHECKPOINT
walks all of the above automatically: it assembles the blocks, runs the
dry run, shows you the output, runs \texttt{--write}, performs the
MASTER-PROMPT.md special case if relevant, and commits. You still see
every step's output --- it is automated, not silent.

\begin{mdframed}[style=warnbox]
\textbf{Only the ai-library-system project may alter MASTER-PROMPT.md.}
A checkpoint from any other project that touches MASTER-PROMPT.md is an
error --- stop and fix the project assignment before writing.
\end{mdframed}

% ---------------------------------------------------------------
\section{NOTE --- Logging Between Checkpoints}

Use NOTE for something worth recording in THREAD.md that does not justify
a full checkpoint --- a decision made mid-session, an observation, a
small correction.

\begin{lstlisting}
echo "Note body." | python3 code/add_note_thread.py [slug] "Topic" --write
\end{lstlisting}

The script appends the note in the correct format and commits THREAD.md
immediately. THREAD.md is never edited by hand for notes --- always
through this script, so the immediate-commit rule is enforced
automatically rather than relying on memory.

% ---------------------------------------------------------------
\section{COMMIT --- Routine Changes}

Use COMMIT for ordinary work that is not part of a version triplet and
not a THREAD.md note --- editing a script in \texttt{code/}, adding a
document to a project's \texttt{docs/} folder, fixing a typo in MAP.md.

\begin{mdframed}[style=stepbox]
\begin{enumerate}
\item Run \texttt{git status} and review what changed.
\item Confirm which files belong in this commit --- COMMIT should not
      silently sweep in unrelated work.
\item Stage exactly those files and commit with a descriptive message.
\end{enumerate}
\end{mdframed}

\begin{lstlisting}
git add [file1] [file2] ...
git commit -m "[descriptive message]"
\end{lstlisting}

% ---------------------------------------------------------------
\section{NEW PROJECT --- Starting Something New}

Use NEW PROJECT to create a new project folder with a persona and a blank
THREAD.md, ready for its first CHECKPOINT.

\begin{mdframed}[style=stepbox]
\begin{enumerate}
\item \textbf{Name the project.} A slug is derived automatically
      (lowercase, hyphens, no punctuation) and confirmed with you.
\item \textbf{Three guided persona questions, asked one at a time:}
      \begin{enumerate}
      \item What is this project, and what is the AI's job in it?
      \item How should the AI behave --- tone, directness, at least one
            hard constraint?
      \item One concrete example exchange: a question you'd ask, and the
            answer you'd want. This shapes behavior more reliably than
            any description.
      \end{enumerate}
\item \textbf{persona.md is generated} from the three answers (ROLE,
      DOMAIN, BEHAVIOUR, EXAMPLES) and shown to you for confirmation
      before anything is written.
\item \textbf{Once confirmed,} the project folder, \texttt{THREAD.md},
      and \texttt{v00} stub triplet are created and committed. The v00
      stubs exist only so RESUME and the pre-commit triplet check both
      work before the first real CHECKPOINT produces v01.
\item \textbf{MAP.md is updated} with a new section for the project.
\end{enumerate}
\end{mdframed}

THREAD.md starts blank apart from its schema. Fill in ``What this project
is'' after the first real working session --- there is no need to
front-load it before any work has happened.

% ---------------------------------------------------------------
\section{Other Navigation Commands}

\begin{center}
\begin{tabular}{@{}lp{8.5cm}@{}}
\toprule
\textbf{Command} & \textbf{What it does} \\
\midrule
SHOW MAP & Summarizes the current state of the library from MAP.md. \\
SHOW THREAD & Summarizes a project's checkpoint log and describes its
arc in plain language. \\
\bottomrule
\end{tabular}
\end{center}

% ---------------------------------------------------------------
\section{What Runs Automatically}
\label{sec:hook-detail}

\subsection{The Pre-Commit Hook}

\texttt{code/githooks/pre-commit} runs on every commit once
\texttt{core.hooksPath} is configured. It has two severities:

\begin{itemize}
\item \textbf{Blocks the commit:} a version triplet is incomplete, or
      triplet version numbers do not match each other (lockstep
      violation).
\item \textbf{Warns but does not block:} a staged file not yet listed
      in MAP.md, a MAP.md entry's path that does not match an actual
      file, THREAD.md entries out of numeric order.
\end{itemize}

\begin{mdframed}[style=notebox]
There is no checksum or tamper-detection step. Git's commit history is
the system of record for integrity --- a separate manifest would only
duplicate what git already verifies on every clone and pull.
\end{mdframed}

\subsection{The ai-library-ops Skill}

A repo-local skill at \texttt{.claude/skills/} (loads automatically in
Claude Code, including cloud and iOS sessions) that drives CHECKPOINT,
NOTE, COMMIT, RESUME, and NEW PROJECT end to end --- the mechanical parts
of every section above. It does not change what any ritual produces; it
only removes the manual copy-paste-and-run steps.

% ---------------------------------------------------------------
\section{Multi-Platform Notes}

\begin{center}
\begin{tabular}{@{}lp{8.5cm}@{}}
\toprule
\textbf{Platform} & \textbf{How MASTER-PROMPT.md loads} \\
\midrule
Claude Code (CLI / cloud / iOS) & Automatically, via \texttt{CLAUDE.md}'s
import on every session start. \\
Cowork & Paste once into the project's system-prompt field; it then
persists for that project. \\
Other AI vendors & No automatic mechanism --- paste MASTER-PROMPT.md at
the start of each session. \\
\bottomrule
\end{tabular}
\end{center}

The rituals themselves (RESUME, CHECKPOINT, NOTE, COMMIT, NEW PROJECT) are
the same regardless of platform --- only how the master prompt gets
loaded, and whether a skill automates the mechanical steps, differs.

% ---------------------------------------------------------------
\section*{Colophon}

\vspace{8pt}
\hrule height 0.4pt
\vspace{12pt}

\begin{center}
{\small\texttt{AI-Library / USER-GUIDE.tex} \quad Version 1.0 \quad \today}\\[6pt]
{\small Plain text. Open format. No vendor. No application required. No expiry.}\\[4pt]
{\footnotesize
This guide explains how to act on the library.\\
MASTER-PROMPT.md is the authority on what the rituals mean.\\
ARCHITECTURE.tex is the authority on why they exist.
}
\end{center}

\end{document}
